\newpage
\begin{center}
	\fontsize{14}{17}\selectfont\bfseries Resumen\addcontentsline{toc}{chapter}{Resumen}
\end{center}
\noindent Esta tesis se refiere al uso del método del paso
fraccionario para la solución numérica de un sistema hiperbólico de
leyes de conservación.
En esta tesis.
Tres aspectos principales del método son considerados.
La primera es la \emph{precisión} y la \emph{eficiencia}
del método de paso fraccionario en relación con los métodos
completamente implícitos.
El segundo tema es la estabilidad de los métodos de paso fraccionario.
En la primera parte, generalizamos el concepto de una solución
clásica y definimos una convergencia en el sentido distribucional
acerca de cuál solución débil corresponde a la solución físicamente
coherente.
También, presentamos métodos populares del esquema de volúmenes
finitos, la convergencia, la consistencia y la estabilidad son
definidos en el contexto vectorial.
En la segunda parte, mostramos los sistemas hiperbólicos a analizar
y desarrollar, con ayuda del programa \textsc{Clawpack} (Conservation
Laws Package)~\cite{Clawpack2025,Mandli2016} y modelamos el flujo de
un fluido de dos fases en un medio poroso.
Finalmente, se discute la calidad de las soluciones aproximadas con
respecto a otro enfoque popular de semidiscretización como el método
de las líneas.

\noindent Palabras clave --- Método de volúmenes finitos, sistema
hiperbólico de leyes de conservación, Clawpack, método de paso
fraccionario.
